\section{Recovering the law from the retained pair data}
\label{p1:close:section}

The stopped expansion requires a positive sum of joint
probabilities.  The component theorem controls centered
coefficients instead.  The two quantities cannot be interchanged
by taking absolute values of the entire expansion: a single good
pair has mass of the same order as the correction we wish to
retain.  We separate that pair explicitly, estimate its failed
extensions after compression, and use the component bound only
for the remaining terms.

Throughout this section, $\mathscr D$ is the class from
Section~\ref{p14:comp:stopped}, $n=|G|$, $\mu=np$, and
\[
 \beta_N=\sum_{e\in E_N}q_e.
\]
All asymptotic constants may depend on the fixed critical band.

\subsection{Two deficits under independent counting weights}

For $r\ge0$ define
\begin{align}
 F_r&=p^r\#\{I\subseteq G:|I|=r,\ I\text{ is not direct}\},
       \label{p1:close:deficit-F}\\
 G_r&=\sum_{e\in E_N}q_ep^{r-2}
      \left\{\binom{n-2}{r-2}-\#\Gamma_e(r-2)\right\}
                   \qquad(r\ge2),
       \label{p1:close:deficit-G}
\end{align}
and put $G_0=G_1=0$.  We use zero for binomial coefficients
whose lower index lies outside the usual range.  In particular,
both sequences vanish for $r>n$, and $F_0=F_1=0$.
The weighted deficits are
\begin{equation}
 F=\sum_{r\ge0}2^rF_r,\qquad
 \mathfrak G=\sum_{r\ge0}2^rG_r.
 \label{p1:close:total-deficits}
\end{equation}
These are deterministic counting deficits.  The added centers
carry independent weights $p$, even when the corresponding
arithmetic events are dependent.

\begin{lemma}[Bounds for the independent deficits]
\label{p1:close:deficit-bounds}
For every fixed $\varepsilon>0$,
\begin{equation}
 F\ll_\varepsilon N^{-4/5+\varepsilon}+B/N,
 \qquad
 \mathfrak G\ll_\varepsilon N^{-1+\varepsilon}.
 \label{p1:close:deficit-rates}
\end{equation}
\end{lemma}

\begin{proof}
Let $\mathcal P$ denote the minimal non-direct families,
including pairs, and set
\[
 S_0=\sum_{C\in\mathcal P}(2p)^{|C|}.
\]
Every non-direct family contains some $C\in\mathcal P$.
After choosing it, sum all added centers with weight $2p$.
The resulting union bound gives
\begin{equation}
 F\le(1+2p)^n S_0\le e^{2\mu}S_0.
 \label{p1:close:ordinary-union}
\end{equation}
For a minimal non-direct family of size at least three,
every proper subfamily is direct.  Expanding the centered
product therefore gives
$c(C)=q(C)-p^{|C|}$, and hence
\[
 p^{|C|}\le q(C)+|c(C)|.
\]
Theorem~\ref{p1:core:primitive}, with fixed weight $2$, bounds
the contribution of these families by
$O_\varepsilon(N^{-4/5+\varepsilon})$.
For pairs, the deterministic host bounds of
Theorem~\ref{p1:arith:pairs}, including the $O(NB)$ contacts
and good-pair hosts, give
\[
 \sum_{\substack{C\in\mathcal P\\|C|=2}}(2p)^2
    \ll_\varepsilon N^{-4/5+\varepsilon}+B/N.
\]
These host counts include incompatible pairs; their independent
weights are still $p^2$.  This proves the bound for $S_0$ and $F$.

Fix now $e\in E_N$ and a family of singleton blocks which is not
direct after adjoining $V_e$.  A minimal non-direct subfamily
either lies entirely among the singletons or contains $V_e$.
The first possibility contains an original member of $\mathcal P$.
Write $(e,C)$ for a minimal non-direct family of the second
kind, where $\varnothing\ne C\subseteq G\setminus e$.
Summing additional centers with their independent weights gives
\begin{equation}
 \mathfrak G
 \le4e^{2\mu}\left\{
       \beta_N S_0+
       \sum_{(e,C)\ {\rm minimal}}
                 2^{|C|}q_ep^{|C|}
             \right\}.
 \label{p1:close:compressed-union}
\end{equation}
The factor $4$ accounts for the two original endpoints of $e$.

To use the compressed estimate, regard
$J_e=J_xJ_y$ as one indicator of mean $q_e$ and set
\[
 \widehat c(e,C)
   =\E\left[(J_e-q_e)\prod_{z\in C}(J_z-p)\right].
\]
For a minimal non-direct family of these blocks, every proper
block subfamily is direct.  The same centered expansion yields
\[
 \widehat c(e,C)=q(e\cup C)-q_ep^{|C|},
 \qquad
 q_ep^{|C|}\le q(e\cup C)+|\widehat c(e,C)|.
\]
Theorem~\ref{p1:core:compressed}, with the pair distinguished,
bounds the second sum in \eqref{p1:close:compressed-union}
by $O_\varepsilon(N^{-1+\varepsilon})$.
The factor $4\,2^{|C|}=2^{|C|+2}$ is exactly the weight of
the original windows; it introduces no loss depending on
the size.

Finally, $\beta_N=O(a_N)$ and
$a_N\asymp LN^{-2/3}$.  The already obtained bound for $S_0$
therefore gives $\beta_NS_0=O_\varepsilon(N^{-1+\varepsilon})$,
after choosing a smaller fixed error margin in the input bounds
if necessary.  Equation \eqref{p1:close:compressed-union}
proves the second assertion.
\end{proof}

The extensions in this proof are summed under the independent
weights.  In particular, no identity
$q(C\cup U)=q(C)p^{|U|}$ has been used for a family which is
not in direct sum.

\subsection{The exact measure retained by the hereditary class}

Let $b=\delta_0+p\Delta$, the Bernoulli law of mean $p$, and
define
\begin{equation}
 \mathsf B_\beta=b^{*n}
          +\beta_N\Delta^{*2}*b^{*(n-2)}.
 \label{p1:close:bernoulli-reference}
\end{equation}
At this stage it is convenient to work with a signed correction
to the binomial law.

\begin{proposition}[Evaluation of the retained polynomial]
\label{p1:close:retained-reference}
For the class $\mathscr D$,
\begin{equation}
 T_{\mathscr D}
 =\mathsf B_\beta-\sum_{r\ge0}(F_r+G_r)\Delta^{*r},
 \qquad
 \|T_{\mathscr D}-\mathsf B_\beta\|_1\le F+\mathfrak G.
 \label{p1:close:retained-identity}
\end{equation}
\end{proposition}

\begin{proof}
The total joint probability of direct families of size $r$ is
$\binom nrp^r-F_r$.
By Lemma~\ref{p1:stop:class}, every non-direct member of
$\mathscr D$ contains a unique good pair.  Equations
\eqref{p1:stop:gamma} and \eqref{p1:stop:direct-pair-probability}
show that their total contribution is
\[
 \sum_{e\in E_N}q_ep^{r-2}\#\Gamma_e(r-2)
   =\beta_N\binom{n-2}{r-2}p^{r-2}-G_r.
\]
This term is zero for $r<2$.  Summing against
$\Delta^{*r}$ and applying the binomial formula gives the
identity.  The norm bound follows from
$F_r,G_r\ge0$ and $\|\Delta^{*r}\|_1=2^r$.
\end{proof}

The coefficient $\beta_N$ is the sum of the joint probabilities
of good pairs.  It will become the covariance correction after
the independent contribution of those pairs and the binomial
diagonal have been accounted for.

\subsection{From centered components to positive boundary mass}

For $K\le L/4$ put
\[
 \mathscr Q=\{I\subseteq G:\rho(I)\le3L/2\},
 \qquad
 R_K=\sum_{\substack{I\in\mathscr Q\setminus\mathscr D\\|I|\le K}}
          2^{|I|}q(I).
\]
The boundary of size at most $K$ is contained in this sum.
Define the component activities
\begin{align}
 A_{\rm good}&=4\sum_{e\in E_N}c(e),\notag\\
 A_{\rm bad}&=4\sum_{\substack{I\subseteq G,\ |I|=2\\I\notin E_N}}
                                      |c(I)|,\notag\\
 B_K&=\sum_{\substack{C\ {\rm connected},\ 3\le|C|\le K\\
                                  \rho_*(C)\le3L/2}}
                              2^{|C|}|c(C)|.
 \label{p1:close:activities}
\end{align}
For a good pair, compatibility and positive nullity give
$c(e)=q_e-p^2>0$.  In particular,
$0\le A_{\rm good}\le4\beta_N=O(a_N)$.

\begin{proposition}[Positive mass outside the retained class]
\label{p1:close:boundary-mass}
For every $K\le L/4$,
\begin{equation}
 R_K\le F+\mathfrak G+
 e^{2\mu}\bigl\{
          \exp(A_{\rm good}+A_{\rm bad}+B_K)-1-A_{\rm good}
          \bigr\}.
 \label{p1:close:boundary-inequality}
\end{equation}
Consequently, uniformly in that range,
\begin{equation}
 R_K\ll_\varepsilon
        N^{-4/5+\varepsilon}+B/N+a_N^2.
 \label{p1:close:boundary-rate}
\end{equation}
\end{proposition}

\begin{proof}
Expand each indicator about its mean:
\begin{equation}
 q(I)=\sum_{J\subseteq I}p^{|I|-|J|}c(J),
 \qquad c(\varnothing)=1.
 \label{p1:close:centered-expansion}
\end{equation}
If $I\in\mathscr Q$, then every $J\subseteq I$ and every
canonical component $C$ of $J$ satisfy
$\rho_*(C)\le\rho(C)\le\rho(I)\le3L/2$.
The first inequality follows from the definition of compatible
nullity, and the second from monotonicity of nullity.

The term $J=\varnothing$, summed over
$I\in\mathscr Q\setminus\mathscr D$, contributes at most $F$:
all such $I$ are non-direct.

For a term $J=e\in E_N$, write $I=e\cup U$.
Since $I\notin\mathscr D$, the family $U$ does not belong to
$\Gamma_e(|U|)$.  Using $0<c(e)\le q_e$ and then dropping
the restrictions on $I$ gives the upper bound
\[
 4\sum_{e\in E_N}q_e
      \sum_{\substack{U\subseteq G\setminus e\\
                          U\notin\Gamma_e(|U|)}}(2p)^{|U|}
       =\mathfrak G.
\]
This is the only part of the argument where a term of order
$a_N$ has to be removed before taking an absolute majorant.

It remains to sum nonempty $J$ which are not a single good pair.
The coefficient $c(J)$ factors over its canonical direct-sum
components.  A singleton component has zero coefficient.
Every remaining nonzero family is therefore a collection of
components of size at least two, containing either a bad pair,
a component of size at least three, or at least two good pairs.
Its component activities are included in
\eqref{p1:close:activities}.

The sum over these collections is bounded by
\begin{equation}
 \sum_{\text{remaining }J}2^{|J|}|c(J)|
 \le \exp(A_{\rm good}+A_{\rm bad}+B_K)-1-A_{\rm good}.
 \label{p1:close:component-exponential}
\end{equation}
Indeed, a collection of $m$ distinct components has $m!$
orders, cancelling the factorial in the exponential series.
Allowing repeated, intersecting or non-independent components
only adds nonnegative terms to the majorant.  The empty
collection and a single good-pair component are precisely
the two terms subtracted on the right.

For each fixed $J$, summing the additional members of
$I\setminus J$ in \eqref{p1:close:centered-expansion}
costs at most
$(1+2p)^{n-|J|}\le e^{2\mu}$.  This proves
\eqref{p1:close:boundary-inequality}.

Theorem~\ref{p1:arith:pairs} and
Theorem~\ref{p1:core:bounded-nullity}, with fixed
$\eta=3/2$ and $w=2$, give
\[
 A_{\rm bad}+B_K
      \ll_\varepsilon N^{-4/5+\varepsilon}+B/N.
\]
Use a small fixed error margin so that this quantity tends to
zero.  For $u,v\ge0$ tending to zero,
$e^{u+v}-1-u=O(u^2+v)$.
The bounds for $A_{\rm good}$, $F$ and $\mathfrak G$
now give \eqref{p1:close:boundary-rate}.
An estimate for a smaller error margin implies the displayed
estimate for every prescribed $\varepsilon>0$.
\end{proof}

\subsection{Converting the binomial reference to Poisson}

Set
\begin{equation}
 C_2=\sum_{\substack{I\subseteq G\\|I|=2}}c(I),
 \qquad
 \alpha_G=\operatorname{Var}(Z_G)-\mu=2C_2-np^2.
 \label{p1:close:variance}
\end{equation}
The second identity is the variance formula for a sum of
indicators of common mean $p$.  Moreover,
\begin{equation}
 \beta_N-C_2
       =|E_N|p^2-\sum_{\substack{|I|=2\\I\notin E_N}}c(I),
 \qquad
 |\beta_N-C_2|
       \le |E_N|p^2+\frac14A_{\rm bad}.
 \label{p1:close:beta-covariance}
\end{equation}
Thus this difference is
$O_\varepsilon(N^{-4/5+\varepsilon}+B/N)$.

\begin{lemma}[Poisson--Charlier conversion]
\label{p1:close:conversion}
With the parameters of the critical band,
\begin{equation}
 \|\mathsf B_\beta-Q_{\mu,\alpha_G}\|_1
 \le 4|\beta_N-C_2|+
           O_C\bigl(N^{-2}+a_N/N\bigr).
 \label{p1:close:conversion-bound}
\end{equation}
\end{lemma}

\begin{proof}
Work in the Banach algebra of absolutely summable signed
measures.  Since $2p<1$ for large $N$, the convolution
logarithm converges absolutely and gives
\[
 \log_*(\delta_0+p\Delta)
     =p\Delta-\frac{p^2}{2}\Delta^{*2}
                     +O_{\ell^1}(p^3).
\]
The convolution exponential of $\mu\Delta$ is the Poisson
probability $\Pi_\mu$, of norm one.  Multiplying the preceding
identity by $n$ and exponentiating therefore yields
\begin{equation}
 b^{*n}
 =\Pi_\mu-\frac{np^2}{2}\Pi_\mu*\Delta^{*2}
           +O_{\ell^1}\bigl(np^3+(np^2)^2\bigr).
 \label{p1:close:binomial-expansion}
\end{equation}
The same calculation, retaining only the leading term, gives
\[
 b^{*(n-2)}
       =\Pi_\mu+O_{\ell^1}(p+np^2).
\]
First replace $\beta_N$ by $C_2$ in
\eqref{p1:close:bernoulli-reference}.
This costs at most $4|\beta_N-C_2|$, because
$b^{*(n-2)}$ is a probability and $\|\Delta^{*2}\|_1=4$.
The two expansions then give coefficient
$C_2-np^2/2=\alpha_G/2$ in front of
$\Pi_\mu*\Delta^{*2}$.  Their remaining error is bounded by
\[
 O_C\bigl(np^3+(np^2)^2+|C_2|(p+np^2)\bigr).
\]
Now $n=O(N)$, $p=O_C(N^{-1})$ and $C_2=O(a_N)$.
This is the error stated in \eqref{p1:close:conversion-bound}.
\end{proof}

\subsection{The first approximation and its uniform boundary}

\begin{proposition}[The first quadratic approximation]
\label{p13:base-approximation}
Uniformly in the critical band, for every \(\varepsilon>0\),
\[
 \|\Law Z_N-Q_{\mu_N,\alpha_N}\|_1
       \ll_{C,\varepsilon}N^{-4/5+\varepsilon}.
\]
The same bound holds with \(\mu_N\) replaced by \(\lambda_N\).
\end{proposition}

This estimate identifies the leading correction and proves the
boundary estimates used below. The optimal scalar reconstruction in
Section~\ref{p13:transfer:section} keeps exact marginals
on the whole set of admissible windows.

\begin{proof}
Choose $\mathscr D$ and $K_N$ as in
Section~\ref{p14:comp:stopped}.  For all sufficiently large $N$,
$K_N\le L/4$.  The stopped identity, the boundary-nullity bound,
Lemma~\ref{p1:stop:tail} and
Proposition~\ref{p1:close:boundary-mass} give
\[
 \|\Law Z_G-T_{\mathscr D}\|_1
 \le R_{K_N}+O_C(N^{-3/2})
 \ll_\varepsilon N^{-4/5+\varepsilon}
                      +B/N+a_N^2+N^{-3/2}.
\]
Proposition~\ref{p1:close:retained-reference} adds the two
deficits $F+\mathfrak G$.
Lemma~\ref{p1:close:conversion}, together with
\eqref{p1:close:beta-covariance}, then gives
\begin{align}
 \|\Law Z_G-Q_{\mu,\alpha_G}\|_1
 &\ll_\varepsilon N^{-4/5+\varepsilon}
             +B/N+a_N^2+N^{-3/2}+N^{-2}+a_N/N \notag\\
 &\ll_\varepsilon N^{-4/5+\varepsilon},
 \label{p1:close:regular-law}
\end{align}
since $B\asymp\log N$ and
$a_N\asymp LN^{-2/3}$.

It remains to restore the excluded windows and the desired
parameters.  The independent reference bounds
\eqref{p1:close:poisson-lipschitz} and
\eqref{p1:close:reference-lipschitz} will transfer the parameters.

Proposition~\ref{p1:arith:core} gives
\[
 \|\Law Z_N-\Law Z_G\|_1
   +|\mu_N-\mu|+|\alpha_N-\alpha_G|
   +|\mu_N-\lambda_N|
      \ll_\varepsilon N^{-1+\varepsilon}.
\]
Apply \eqref{p1:close:reference-lipschitz} first to transfer
$(\mu,\alpha_G)$ to $(\mu_N,\alpha_N)$, and then to replace
$\mu_N$ by $\lambda_N$ while retaining $\alpha_N$.
Both costs are absorbed in \eqref{p1:close:regular-law}.
This proves the two asserted forms of the approximation,
uniformly in the fixed critical band.

Finally, if $0<\varepsilon<2/15$, then
\[
 \frac{N^{-4/5+\varepsilon}}{a_N}
       =O_C\left(\frac{N^{-2/15+\varepsilon}}{\log N}\right)
       \longrightarrow0.
\]
Thus the corrected error is smaller than the retained pair term.
\end{proof}


\subsection{The total-variation coefficient}
\label{p14:comp:tv}

\begin{proof}[Proof of Corollary~\ref{p1:cor:tv}]
Theorem~\ref{p1:thm:main} and the reverse triangle
inequality first give
\[
 \TV(\Law Z_N,\Pi_{\mu_N})
 =\alpha_N\mathcal C(\mu_N)
            +O_{C,\varepsilon}(N^{-4/3+\varepsilon}).
\]
Use \(|\mu_N-\lambda_N|\ll N^{-1+\varepsilon}\) and
\eqref{p1:close:reference-lipschitz} to obtain
\[
 \Law Z_N-\Pi_{\lambda_N}
   =\frac{\alpha_N}{2}\Pi_{\lambda_N}*\Delta^{*2}
          +O_{\ell^1,C,\varepsilon}(N^{-1+\varepsilon}).
\]
The reverse triangle inequality, followed by the factor
$1/2$ converting the $\ell^1$ norm to total variation,
yields
\begin{equation}
 \TV(\Law Z_N,\Pi_{\lambda_N})
      =\alpha_N\mathcal C(\lambda_N)
          +O_{C,\varepsilon}(N^{-1+\varepsilon}).
 \label{p1:close:tv-exact-parameter}
\end{equation}
The pair asymptotic from
Theorem~\ref{p1:arith:pairs} now gives
\begin{equation}
 \TV(\Law Z_N,\Pi_{\lambda_N})
 \sim \mathcal C(\lambda_N)
       \frac{r_N+2}{18\,2^{r_N/3}}\,
       \lambda_N^{5/3}L N^{-2/3},
 \qquad r_N=L\bmod3.
 \label{p1:close:tv-asymptotic}
\end{equation}
The lower bound in \eqref{p1:close:constant-bounds} is uniform
on the fixed critical band, and the remainder is
$o(LN^{-2/3})$.  This proves the sharp order asserted in the
corollary.  The equivalence is relative to the displayed
oscillating expression; no convergence of its phase or
intensity factors is required.
\end{proof}

Replacing $\alpha_N$ by its leading equivalent proves
\eqref{p1:close:tv-asymptotic}.  It does not by itself preserve
the error $O(N^{-1+\varepsilon})$ in
\eqref{p1:close:tv-exact-parameter}, which uses the exact
second factorial cumulant.

For the representative case $\lambda=1$, the only negative
masses of $\Pi_1*\Delta^{*2}$ occur at $k=1$ and $k=2$.
They are $-e^{-1}$ and $-e^{-1}/2$, respectively.
The total mass is zero, so
\[
 \|\Pi_1*\Delta^{*2}\|_1=\frac3e,
 \qquad
 \mathcal C(1)=\frac3{4e}.
\]
Along $N=2^L$ with $L$ divisible by three,
\begin{equation}
 \TV(\Law Z_N,\Pi_1)
       \sim \frac1{12e}\,L N^{-2/3}.
 \label{p1:close:dyadic-example}
\end{equation}
This computes the coefficient in the example announced in
the introduction.

The total-variation asymptotic uses the exact-marginal reconstruction
and the explicit pair coefficient. The independent positivity and
reference-continuity estimates do not use the reconstruction.
